\section{Introduction}

Inductor, resistor, capacitor (LRC) circuits are a fundamental configuration of components that allows us to examine
the behaviour between electricity and magnetism.
This is because capacitors store energy in an electric field, and inductors store energy in its generated magnetic
field from the change in current, and a discharged circuit demonstrates the energy transfer between them.
Understanding LRC circuits is fundamental to modern electronics, and is responsible for radio
frequency communications from integrated
antennae in PCBs, motor drivers, and capacitive touch sensors, for example.

From Kirchoff's voltage law, the primary differential equation that models the discharge behaviour is

\begin{equation}\label{eq:diffyq}
L \odv[2]{i}{t} + R \odv{i}{t} + \frac{i}{C} = 0.
\end{equation}

Where $L$, $R$, $C$, $i$ are the inductance, system resistance, capacitance, and current, respectively.
The general solutions to this equation can be reduced to three configurations that display differing behaviour by
selecting
values for
$R$:
Normal ,
critical, and
oscillatory
damping; Eqs.~\eqref{eq:norm_damp}, ~\eqref{eq:crit_damp}, and \eqref{eq:osc_damp}, respectively.

\begin{equation}\label{eq:norm_damp}
\left( \frac{R}{2 L}\right)^2 - \frac{1}{L C} > 0
\end{equation}

\begin{equation}\label{eq:crit_damp}
\left( \frac{R}{2 L}\right)^2 - \frac{1}{L C} = 0
\end{equation}

\begin{equation}\label{eq:osc_damp}
\left( \frac{R}{2 L}\right)^2 - \frac{1}{L C} < 0
\end{equation}

For the critically damped case, theoretical predictions can be made on the voltage signals.
The theoretical equations for current and voltage are given by the following:

\begin{equation}
i = \frac{V_0}{L} t e^{- \frac{R}{2 L}t}.\label{eq:i_theo}
\end{equation}

\begin{equation}\label{eq:vl_theo}
V_L = -V_0 \left(1 - \frac{R}{2 L} t\right) e^{- \frac{R}{2 L}t}.
\end{equation}

\begin{equation}\label{eq:vc_theo}
V_C = V_0 \left(1 + \frac{2}{R C} t\right) e^{- \frac{R}{2 L}t}.
\end{equation}

$V_0$, $V_L$, $V_C$ are the values for the steady state voltage, inductor, capacitor electric potentials,
respectively.
For the oscillatory case, the theoretical frequency of oscillation is given by

\begin{equation}\label{eq:freq}
j \omega = \sqrt{\frac{R^2}{4 L^2} - \frac{1}{L C}}.
\end{equation}

The energetics of the system can be calculated using potential energy equations:

\begin{equation}\label{eq:el}
    E_L = \frac{1}{2} Li^2
\end{equation}

\begin{equation}\label{eq:ec}
    E_C = \frac{1}{2} C V_C^2
\end{equation}

And the thermal heat energy dissipated is calculated using

\begin{equation}\label{eq:heat}
    E_H = \int_{t_n}^{t_{n+\Delta t}} Ri^2 \dif t \approx R i_{AV}^2 \Delta t,
\end{equation}

where $\Delta t$ is the evaluated time step between samples, and $i_{AV}$ is the average current over the time
interval given by

\begin{equation}\label{eq:i_av}
    i_{AV} = \frac{V_R}{R_{res}}.
\end{equation}

